\subsection{The uniform cuspidal sector}

\begin{statement}{P4.43}{Proved}
In the requested full-fibration classes, a support with
positive-genus ruled resolution cannot have positive
Euler number and generic transverse type $A_2$ along
every component of its critical curve.
\end{statement}
\begin{proof}
We use the exact normalization map
$\nu:N\longrightarrow S$ and its minimal resolution
$\rho:T\longrightarrow N$. A germ of $S$ cannot have
two distinct analytic hypersurface branches. Indeed,
their intersection has dimension one in the smooth
ambient threefold. A general point of an irreducible
component of that intersection has at least two local
hypersurface branches, so its transverse curve germ is
reducible. The transverse cusp $A_2$ is irreducible.
This contradicts the stated uniform transverse type.
The points of the finite normalization over a germ are
its analytic branches. Consequently $\nu$ is bijective.
It is proper, so its continuous map from a compact space
to the Hausdorff space $S$ is a homeomorphism.

Let $B$ be the base curve of the minimal ruled surface,
of genus $g\ge1$, and let $H$ be the horizontal section
of coefficient two in the total pole divisor. Point
blowups preserve the fundamental group, and the ruling
gives $\pi_1(T)=\pi_1(B)$. By the fibre-tree calculation
in P4R.7, every vertical exceptional configuration is a
tree of rational curves. Collapsing any such connected
tree preserves the fundamental group: apply van Kampen
to its tubular neighbourhood and the complement, or
equivalently to the cofibration quotient by that
simply connected tree. These quotients are the actual
topological contractions of the resolution map.

If $H$ is not contracted by $\rho$, all resolution
exceptional configurations are of this vertical kind.
It follows that $\pi_1(N)=\pi_1(B)$. If $H$ is
contracted, its connected exceptional configuration
contains $H$. The inclusion $H\hookrightarrow T$
induces an isomorphism on fundamental groups, since
the ruling restricts to the identity section and all
point blowups preserve that isomorphism. The quotient
which collapses this configuration therefore kills
the whole group. The other contracted configurations
are rational trees, giving $\pi_1(N)=1$.

On the other hand, the proved positive ruled reduction,
including P4R.13, gives a conductor curve $C$ with two
rational normalized components and one graph cycle.
The conductor cofibration pushout in P4R.7 gives
$\pi_1(S)=\pi_1(C)=\mathbb Z$. This contradicts the
homeomorphism $N\simeq S$: neither the trivial group
nor the genus-$g$ surface group is $\mathbb Z$, as their
abelianizations have ranks zero and $2g$, respectively.
\end{proof}
